\documentclass[10pt]{article}
\usepackage[margin=20mm]{geometry}
\usepackage{amsmath,amssymb,fontspec,unicode-math,xurl,hyperref}
\IfFontExistsTF{DejaVu Serif}{\setmainfont{DejaVu Serif}}{\setmainfont{Latin Modern Roman}}
\setmathfont{latinmodern-math.otf}
\setlength{\parindent}{0pt}\setlength{\parskip}{6pt}\sloppy
\begin{document}
Abstract: This report checks posterior normalization, the balanced-transport marginal identity, the published ACAUC pseudocode, and the proposed fallback under distribution shift. Equation (5) and the finite empirical marginal-calibration identity are correct under their stated normalization. Two appendix algorithms omit essential normalization factors. A two-point entropic-transport calculation also shows that separated representation supports do not force uniform coupling. These are bounded specification and extrapolation defects; the reported empirical performance was not reproduced. Scientific adjudication remains pending.\par
Author: AI-assisted technical review\par
Date:\par
Title: Technical Verification of RoPE\par
\# Identity and scope\par
The reviewed paper is Antoine Wehenkel, Juan L. Gamella, Ozan Sener, Jens Behrmann, Guillermo Sapiro, Jörn-Henrik Jacobsen, and Marco Cuturi, Addressing Misspecification in Simulation-based Inference through Data-driven Calibration, ICML 2025, PMLR 267, pp. 65949–65980 [Retained latex: \textbackslash{}cite\{paper\}]; local identifier 0187/I25-43556. The source is the 32-page proceedings PDF with embedded appendices. Page references are one-based PDF pages, matching printed labels.\par
The check covered Sections 2–3, the metric and interpretation passages in Sections 5–6, Algorithm 1, Appendices E–H, and Algorithm 9. Selected experimental settings and figures were read for scope. A matching RoPE implementation was not available in the staged sources. No author code, datasets, model weights, or experiments were executed or downloaded. The causal-chamber repository mentioned by the paper is not treated as a RoPE implementation. Separate OpenReview attachment access remains unresolved.\par
\# Correct core identities and their limits\par
For \(n_o\) observations and \(n_s\) simulations, the transport matrix satisfies \[\sum_jP_{ij}=1/n_o.\] Thus Equation (5)’s weights \(\alpha_{ij}=n_oP_{ij}\) are nonnegative and sum to one. The mixture of normalized simulation posteriors is consequently a normalized posterior.\par
Equation reference: eq:marginal. For balanced transport, \(\sum_iP_{ij}=1/n_s\). Direct summation gives the exact empirical identity \[
\frac1{n_o}\sum_i\widetilde p(\theta\mid x_o^i)
=\frac1{n_s}\sum_j\widetilde p(\theta\mid x_s^j).\] Under the assumed simulator-marginal calibration and the simulation-sampling limit, the right side approaches the prior. Appendix F appropriately calls this a necessary marginal-calibration property; it is not a guarantee of conditional calibration, joint credible-region coverage, or correct coupling. The paper also explicitly acknowledges that ACAUC does not assess joint calibration. No defect is established in ([eq:marginal] (\#eq:marginal)) under these conditions.\par
\# Substantive findings\par
Problem 1 (Algorithm 1 produces a sub-probability density as printed).\par
\#\# Claim, assumptions, and calculation\par
Algorithm 1, Step 5, p. 17, uses \(\sum_jP_{ij}\widetilde p(\theta\mid x_s^j)\) as its posterior, although its Step 4 retains the row constraint defining \(\mathcal B_o\) on p. 3. For normalized components and any \(n_o>1\), \[\int\sum_jP_{ij}\widetilde p(\theta\mid x_s^j)\,d\theta
=\sum_jP_{ij}=1/n_o\ne1.\] For the explicit admissible balanced matrix \(P_{ij}=1/16\) with \(n_o=n_s=4\), the integral is \(1/4\). Multiplying by \(n_o\) gives one. If interpreted as a log-density, the unnormalized quantity differs by \(-\log n_o\); this is a conditional implication, not an observed experimental offset.\par
\#\# Counterevidence, impact, and correction\par
Main Equation (5) and Appendix F both include the missing factor. They supply a precise, analytically verified repair: use \(\alpha_{ij}=n_oP_{ij}\). The intact core identity is important counterevidence against attributing this omission to the experiments. The supported grade is Minor for inconsistent executable specification, with historical implementation impact Needs Evidence. The grade is not based on the appendix location; it reflects that the main mathematical definition and its proof already give the normalized algorithm.\par
Problem 2 (Algorithm 9’s unnormalized rank sum does not estimate an average calibration area).\par
\#\# Claim and source\par
Section 5, p. 7, defines ACAUC as an average over observations, parameter dimensions, and credible levels. Appendix J, Algorithm 9, p. 28, instead sums ordered rank-derived levels minus \(i/N\) at line 16, then divides only by the parameter count \(K\). The original page image confirms the absence of a sample-count normalization. There are also inconsistent \(M/N\) names and list/scalar assignments; the check below grants those elementary implementation repairs and gives the algorithm its most favorable coherent interpretation.\par
\#\# Complete analytic check\par
Take one parameter, continuous correctly specified posteriors, independent observations, and exact posterior ranks (the infinite posterior-draw limit). Then the probability integral transforms \(U_i\) are independent uniform variables. Interpret \(N\) as the number of evaluated observations and line 16 as \[A_N=\sum_{i=1}^{N}\left(U_{(i)}-\frac{i}{N}\right).\] Sorting preserves the sum. Therefore \[\mathbb E A_N=-\frac12,\qquad
\operatorname{Var}(A_N)=\frac{N}{12}.\] The estimator neither converges to zero for a calibrated model nor has the scale of an average area. The concrete grid \(U_i=(i-\frac12)/N\) gives \(A_N=-1/2\) for every \(N\). With \(N=4\), the independent exact-fraction calculation obtains \(-1/2\) rather than a normalized average.\par
\#\# Counterevidence and bounded repair\par
The main metric definition includes the correct observation averaging. Its displayed scores and calibration curves cannot be assumed to come from the literal pseudocode, especially because the latter also contains implementation ambiguities. No original implementation or metric arrays were available, so the empirical calibration claims remain unverified rather than disproved by this finding.\par
The exact connection to the main metric requires an interval convention. For any nested credible-level family \(C_i(\alpha)\), set \[V_i=\inf\bigl(\{\alpha\in[0,1]:\theta_i\in C_i(\alpha)\}\cup\{1\}\bigr).\] Apart from endpoints of zero integration measure, membership is \(\mathbf{1}\{V_i\le\alpha\}\). The single-coordinate empirical signed area in the main definition is therefore exactly \[\frac1N\sum_i\int_0^1
 [\alpha-\mathbf{1}\{\theta_i\in C_i(\alpha)\}]\,d\alpha
 =\frac1N\sum_i V_i-\frac12.\] For one-sided lower credible intervals, \(V_i=U_i=F_i(\theta_i)\); for central equal-tailed intervals, \(V_i=|2U_i-1|\). Both are uniform under a correct continuous posterior, so the normalization diagnosis survives either convention. Under misspecification their distributions differ, and the intended interval definition must be stated.\par
For the corresponding ordered \(V_i\), dividing the right-endpoint rank sum by \(N\) removes the divergent scale; replacing \(i/N\) by \((i-\frac12)/N\) then gives exactly \(\overline V-1/2\). Under independent correct posteriors, the right-endpoint version has mean \(-1/(2N)\) and variance \(1/(12N)\). For \(N=4\) with every \(V_i=0.99\), the literal sum is \(1.46\), outside the signed area’s range \([-1/2,1/2]\), while the exact area is \(0.49\). The metric can also cancel positive and negative calibration errors; it is not alone a complete calibration certificate. The supported defect is Minor, with historical metric implementation and interval semantics Needs Evidence.\par
Problem 3 (Disjoint support does not ensure an automatic uniform-transport fallback).\par
\#\# Claim and assumptions\par
Section 6, p. 9, discusses disjoint calibration/test support and states that a lack of meaningful matches will make the transport uniform and return the prior. This extrapolation is outside the principal i.i.d. calibration setting. The issue is whether the transport objective itself supplies the claimed detection mechanism. Section 3.1 specifies an unsquared Euclidean feature distance; the counterexample uses exactly that distance.\par
\#\# An admissible transport calculation\par
Let the two simulation representations be \((-R,0),(R,0)\) and two test representations be \((-R,R),(R,R)\), with uniform marginals, \(R>0\), balanced transport, and entropy weight \(\gamma>0\). The supports are disjoint, with minimum distance \(R\). The Euclidean cost matrix is \[C=R\begin{pmatrix}1&\sqrt5\\\sqrt5&1\end{pmatrix}.\] Every feasible coupling with these two uniform marginals has the form \[P=\frac12\begin{pmatrix}q&1-q\\1-q&q\end{pmatrix}.\] Strict convexity of entropy ensures uniqueness. Differentiating the objective in \(q\) yields \[-R(\sqrt5-1)+\gamma\log\frac q{1-q}=0,\qquad
q=\frac1{1+\exp[-R(\sqrt5-1)/\gamma]}.\] For \(R=10,\gamma=0.5\), \(q=0.9999999999816493\). Uniform conditional weights would be \(1/2\). As the supports become farther apart, \(q\) tends to one, not to \(1/2\). The assumptions impose no feature-norm bound or constraint on the out-of-distribution extension of the learned representation that excludes this configuration. Semantic correctness of the near-diagonal matching is not determined by its distance ranking. To connect this geometry to the paper’s calibration/test claim, let the calibration representation lie on the same lower line as the simulation representation. It can match the parameter perfectly there. A test distribution on the disjoint upper line may reverse the parameter-to-position association while the representation map retains the displayed geometric positions. Fine-tuning on the lower calibration support constrains neither that out-of-support association nor the relative test costs. The same objective then yields a confident but semantically wrong match. This is an admissible finite transport illustration of the claimed fallback mechanism; it is not a reproduced neural-density experiment.\par
\#\# Scope and correction\par
This disproves automatic fallback from separation or poor generalization alone; it does not disprove the paper’s observed light-tunnel results. Large \(\gamma\) relative to cost differences, or essentially equal transport-relevant costs, can produce a uniform coupling. Increasing all distances is insufficient. The supported correction is to make fallback conditional on the cost geometry/regularization, or to provide a separately validated mismatch detector. No such detector is invented here. This is Minor for the unsupported extrapolation beyond the evaluated conditions, with no demonstrated failure of the principal i.i.d. results.\par
\# Overall impact and unexamined work\par
The highest supported grade is Minor. Main Equation (5), its normalization, and the balanced empirical marginal identity survive. Algorithmic normalization and metric correspondence require correction or clarification; actual experimental provenance is unavailable. The distribution-shift example bounds a discussion claim rather than asserting failure on the published benchmarks.\par
Appendix G explicitly states that unrestricted NPE optimality cannot establish minimality of its representation. Its information-bottleneck proposal is therefore not treated here as a proved minimal-sufficiency theorem. Experimental LPP, full calibration curves, hyperparameter choices, neural training, and raw data splits were not independently reproduced. Later revisions and errata were not exhaustively checked. Scientific adjudication remains pending; human reviews are empty. The issue grades below are unchanged. Adoption has not been requested or verified. For P2, \(V=U\) for lower one-sided intervals and \(V=|2U-1|\) for equal-tailed central intervals; the interval convention is part of the unresolved implementation binding.\par
Issue assessment | Repair verification | Adoption | Remaining work\par
Issue assessment | Repair verification | Adoption | Remaining work\par
P1: Minor; appendix mixture omits a conditional-mass factor. | Multiplication by \(n_o\) restores normalization analytically. | Not requested. | Bind the historical implementation; main Equation 5 remains correct.\par
P2: Minor; appendix rank-area estimator lacks consistent normalization. | Exact signed area \(\overline V-1/2\) checked for the stated nested-interval entry variable. | Not requested. | Identify interval semantics, normalization and original metric arrays.\par
P3: Minor; separation alone does not imply uniform entropic transport. | Exact unsquared-cost counterexample checked; fallback conditions clarified. | Not requested. | Specify suitable cost geometry/regularization and validate any proposed fallback detector.\par
9 A. Wehenkel et al. Addressing Misspecification in Simulation-based Inference through Data-driven Calibration. PMLR 267 (2025), 65949–65980. https://proceedings.mlr.press/v267/wehenkel25a.html. Reviewed proceedings PDF.\par
\# Evidence and actual work\par
Original PDF SHA-256:\par
> .\par
The directory `verification/0187` contains the source manifest binding, extracted text, source-page images 4, 9, 17, and 28 actually inspected by the parent, and independent scalar checks with executed results. The scripts use only standard-library arithmetic and do not run author code. A preliminary squared-distance diagnostic was superseded after checking the original norm notation; its files are retained and are explicitly excluded from the final transport argument. The final check uses the unsquared norm.\par
All scientific integration was performed by the parent AI reviewer. Native auxiliary separately challenged the three findings, rechecked the original source images and independently computed the normalization, interval-entry area and unsquared-cost transport examples. The parent read the challenge and incorporated its exact interval convention and calibration-to-test explanation. Its executed checks and source bindings are preserved separately; no additional parent execution of that challenge script is claimed. No human review or adoption is asserted. Exact TeX/PDF build status, report visual-QA status, execution records, paths, and evidence hashes are stored in the associated provenance file.\par
\end{document}
